\newpage
\chapter*{\centering Abstract}
\textit{
Underwater object detection is a critical capability for autonomous marine robots, enabling applications such as species monitoring, underwater archaeology, and ocean resource management. The URPC (Underwater Robot Picking Contest) benchmark presents a challenging four-class detection problem — holothurian, echinus, scallop, and starfish — in turbid, low-contrast, and color-degraded underwater imagery.

This work presents \textbf{HTDet}, a Hybrid Transformer-based detector built on the MobileViT backbone and a RetinaNet-style detection head, designed explicitly for underwater scenes and FPGA deployment. The core architecture pairs a MobileViT-S backbone (5.58 M parameters) with a Feature Pyramid Network (FPN) neck, achieving a baseline mAP$_{50}$ of \textbf{75.5\%} on the URPC validation set. We further propose a \textbf{Frequency-Aware FPN (FA-FPN)} neck, which exploits the frequency-selective nature of underwater image degradation by learning per-channel spectral gates and frequency-band decompositions at each pyramid level. FA-FPN raises the best mAP$_{50}$ to \textbf{77.6\%}, a +2.1\% absolute gain over the standard FPN neck.

To enable real-time inference on edge hardware, we develop a complete FPGA deployment pipeline. We apply structured channel pruning (30\% width reduction), reducing model parameters by 52\% while retaining an mAP$_{50}$ of \textbf{75.3\%}. For a compact 192$\times$192 inference model, knowledge distillation from the full model recovers accuracy to \textbf{71.8\%} mAP$_{50}$. Post-Training Quantization (PTQ) with W8A32 (INT8 weights, FP32 activations) produces an average SQNR of \textbf{43.8 dB} and cosine similarity $\approx$1.000 across all 53 convolutional layers, validating near-lossless quantization. A complete C/HLS implementation of the 192-resolution model — covering the backbone, FPN neck, and RetinaNet head — is validated against the PyTorch reference, with C-Simulation (CSIM) results matching float32 Python detections on all test images.
}
